% ============================================================
%  Chapter 8: Charge as a Boundary Property
%  Source: bounded-phase-space-trajectory.tex, Consequences 26--27
% ============================================================

\epigraph{%
  Electric charge is not a substance in the interior.\\
  It is a topological fact about the boundary.%
}{}

\section{Charge Emergence from Partition Boundaries}
\label{sec:ch8-emergence}

Electric charge is usually taken as a primitive property of matter, defined
by its coupling to the electromagnetic field.
In the partition framework, charge is not primitive: it is the topological
invariant of a partition boundary.

\begin{definition}[Partition boundary]
  The \emph{partition boundary} $\partial A$ of a cell $A \in \Partition_n$
  is the zero-measure set separating $A$ from adjacent cells.
  It carries the information about the transition between distinct partition
  cells.
\end{definition}

\begin{consequence}{Charge Emergence}{chargeemergence}
  Electric charge is a property of partition boundaries, not of partition
  interiors.
  Specifically, the total charge $q$ enclosed in a region $V$ equals the
  total partition-boundary content of $\partial V$:
  \begin{equation}
    q = \oint_{\partial V} \nabla\Depth \cdot d\mathbf{A},
  \end{equation}
  where $\Depth$ is the partition depth field.
\end{consequence}

\begin{proof}[Proof sketch]
  A categorical observation (Consequence~\ref{cons:categorical}) can only
  detect partition boundaries: a finite-information observer identifies
  which cell they are in, and the boundary is the transition between cells.
  The divergence-free condition on the partition depth gradient in the bulk
  (no sources) combined with Gauss's theorem produces the surface-integral
  form of charge.
  The electromagnetic Gauss law $\nabla \cdot \mathbf{E} = \rho/\varepsilon_0$
  is the continuum limit of this partition-boundary condition.
\end{proof}

\begin{remark}[Five independent proofs]
  The source paper~\cite{Sachikonye2024BPS} provides five independent proofs
  of charge as a boundary property:
  \begin{enumerate}
    \item Via the categorical observation argument (given above).
    \item Via the topological invariant of the partition boundary
      (holonomy of the Fisher-metric connection).
    \item Via dimensional reduction: 3D bulk compresses to 2D boundary
      (analogue of 't Hooft holography).
    \item Via the Loschmidt argument: only boundary crossings are
      irreversible events.
    \item Via the Noether argument: the symmetry under reparametrisation of
      cell labels within a shell generates a conserved boundary current,
      which is identified as electric charge.
  \end{enumerate}
  All five proofs give identical quantisation conditions.
\end{remark}

\section{Charge Quantisation}
\label{sec:ch8-quantisation}

\begin{consequence}{Charge Quantisation}{chargequantisation}
  Electric charge is quantised in units of a fundamental charge $e$:
  \begin{equation}
    q \in \{0, \pm e, \pm 2e, \ldots\}.
  \end{equation}
\end{consequence}

\begin{proof}[Proof sketch]
  The partition boundary $\partial A$ is a measurable set with integer
  topological invariant (the winding number of the boundary map).
  The charge attached to a boundary cell is this winding number times the
  fundamental charge unit.
  The winding number is integer-valued because $\pi_1(SO(3)) = \mathbb{Z}/2$
  and each full boundary crossing increments the winding number by one.
  Fractional charges (as in quarks) correspond to partial boundary crossings
  that are path-opaque (virtual substates, Chapter~\ref{ch:subnucleon}).
\end{proof}

\section{The Coulomb Field as Partition Depth Gradient}
\label{sec:ch8-coulomb}

The electric field outside a charged partition boundary is the gradient of the
partition depth field:
\begin{equation}
  \mathbf{E}(\mathbf{r}) = -\nabla\Depth(\mathbf{r}).
\end{equation}
The divergence-free condition in the bulk ($\nabla^2\Depth = 0$ away from
boundaries) has the unique spherically symmetric solution
$\Depth \propto 1/r$, giving the Coulomb field $\mathbf{E} \propto 1/r^2$.

\begin{consequence}{Coulomb Envelope}{coulombenvelope}
  The partition depth field outside a charge source falls off as $1/r$,
  and the partition depth gradient (electric field) falls off as $1/r^2$,
  in exact agreement with Coulomb's law.
\end{consequence}

\begin{remark}
  The electromagnetic field is not a separate physical entity in this
  framework.
  It is the spatial extension of the partition structure: the gradient of the
  partition depth field in regions accessible to a bounded chain of partition
  boundaries.
  This is compatible with all electromagnetic observations but removes the
  conceptual need for a ``field as substance''.
\end{remark}

\bigskip
\begin{tcolorbox}[colback=crownblue!5, colframe=crownblue!60!black,
                  title={\textbf{Chapter 8 Summary}}]
\begin{itemize}[noitemsep]
  \item Electric charge is a topological invariant of partition boundaries,
    not a property of partition interiors.
  \item Charge quantisation follows from the integer winding number of the
    boundary map.
  \item The Coulomb field is the gradient of the partition depth field in
    the bulk; its $1/r$ form follows from the divergence-free condition.
  \item Maxwell's Gauss law is the continuum limit of the partition-boundary
    charge condition.
\end{itemize}
\end{tcolorbox}
